\documentclass[11pt]{article}

    \usepackage[breakable]{tcolorbox}
    \usepackage{parskip} % Stop auto-indenting (to mimic markdown behaviour)
    

    % Basic figure setup, for now with no caption control since it's done
    % automatically by Pandoc (which extracts ![](path) syntax from Markdown).
    \usepackage{graphicx}
    % Keep aspect ratio if custom image width or height is specified
    \setkeys{Gin}{keepaspectratio}
    % Maintain compatibility with old templates. Remove in nbconvert 6.0
    \let\Oldincludegraphics\includegraphics
    % Ensure that by default, figures have no caption (until we provide a
    % proper Figure object with a Caption API and a way to capture that
    % in the conversion process - todo).
    \usepackage{caption}
    \DeclareCaptionFormat{nocaption}{}
    \captionsetup{format=nocaption,aboveskip=0pt,belowskip=0pt}

    \usepackage{float}
    \floatplacement{figure}{H} % forces figures to be placed at the correct location
    \usepackage{xcolor} % Allow colors to be defined
    \usepackage{enumerate} % Needed for markdown enumerations to work
    \usepackage{geometry} % Used to adjust the document margins
    \usepackage{amsmath} % Equations
    \usepackage{amssymb} % Equations
    \usepackage{textcomp} % defines textquotesingle
    % Hack from http://tex.stackexchange.com/a/47451/13684:
    \AtBeginDocument{%
        \def\PYZsq{\textquotesingle}% Upright quotes in Pygmentized code
    }
    \usepackage{upquote} % Upright quotes for verbatim code
    \usepackage{eurosym} % defines \euro

    \usepackage{iftex}
    \ifPDFTeX
        \usepackage[T1]{fontenc}
        \IfFileExists{alphabeta.sty}{
              \usepackage{alphabeta}
          }{
              \usepackage[mathletters]{ucs}
              \usepackage[utf8x]{inputenc}
          }
    \else
        \usepackage{fontspec}
        \usepackage{unicode-math}
    \fi

    \usepackage{fancyvrb} % verbatim replacement that allows latex
    \usepackage{grffile} % extends the file name processing of package graphics
                         % to support a larger range
    \makeatletter % fix for old versions of grffile with XeLaTeX
    \@ifpackagelater{grffile}{2019/11/01}
    {
      % Do nothing on new versions
    }
    {
      \def\Gread@@xetex#1{%
        \IfFileExists{"\Gin@base".bb}%
        {\Gread@eps{\Gin@base.bb}}%
        {\Gread@@xetex@aux#1}%
      }
    }
    \makeatother
    \usepackage[Export]{adjustbox} % Used to constrain images to a maximum size
    \adjustboxset{max size={0.9\linewidth}{0.9\paperheight}}

    % The hyperref package gives us a pdf with properly built
    % internal navigation ('pdf bookmarks' for the table of contents,
    % internal cross-reference links, web links for URLs, etc.)
    \usepackage{hyperref}
    % The default LaTeX title has an obnoxious amount of whitespace. By default,
    % titling removes some of it. It also provides customization options.
    \usepackage{titling}
    \usepackage{longtable} % longtable support required by pandoc >1.10
    \usepackage{booktabs}  % table support for pandoc > 1.12.2
    \usepackage{array}     % table support for pandoc >= 2.11.3
    \usepackage{calc}      % table minipage width calculation for pandoc >= 2.11.1
    \usepackage[inline]{enumitem} % IRkernel/repr support (it uses the enumerate* environment)
    \usepackage[normalem]{ulem} % ulem is needed to support strikethroughs (\sout)
                                % normalem makes italics be italics, not underlines
    \usepackage{soul}      % strikethrough (\st) support for pandoc >= 3.0.0
    \usepackage{mathrsfs}
    

    
    % Colors for the hyperref package
    \definecolor{urlcolor}{rgb}{0,.145,.698}
    \definecolor{linkcolor}{rgb}{.71,0.21,0.01}
    \definecolor{citecolor}{rgb}{.12,.54,.11}

    % ANSI colors
    \definecolor{ansi-black}{HTML}{3E424D}
    \definecolor{ansi-black-intense}{HTML}{282C36}
    \definecolor{ansi-red}{HTML}{E75C58}
    \definecolor{ansi-red-intense}{HTML}{B22B31}
    \definecolor{ansi-green}{HTML}{00A250}
    \definecolor{ansi-green-intense}{HTML}{007427}
    \definecolor{ansi-yellow}{HTML}{DDB62B}
    \definecolor{ansi-yellow-intense}{HTML}{B27D12}
    \definecolor{ansi-blue}{HTML}{208FFB}
    \definecolor{ansi-blue-intense}{HTML}{0065CA}
    \definecolor{ansi-magenta}{HTML}{D160C4}
    \definecolor{ansi-magenta-intense}{HTML}{A03196}
    \definecolor{ansi-cyan}{HTML}{60C6C8}
    \definecolor{ansi-cyan-intense}{HTML}{258F8F}
    \definecolor{ansi-white}{HTML}{C5C1B4}
    \definecolor{ansi-white-intense}{HTML}{A1A6B2}
    \definecolor{ansi-default-inverse-fg}{HTML}{FFFFFF}
    \definecolor{ansi-default-inverse-bg}{HTML}{000000}

    % common color for the border for error outputs.
    \definecolor{outerrorbackground}{HTML}{FFDFDF}

    % commands and environments needed by pandoc snippets
    % extracted from the output of `pandoc -s`
    \providecommand{\tightlist}{%
      \setlength{\itemsep}{0pt}\setlength{\parskip}{0pt}}
    \DefineVerbatimEnvironment{Highlighting}{Verbatim}{commandchars=\\\{\}}
    % Add ',fontsize=\small' for more characters per line
    \newenvironment{Shaded}{}{}
    \newcommand{\KeywordTok}[1]{\textcolor[rgb]{0.00,0.44,0.13}{\textbf{{#1}}}}
    \newcommand{\DataTypeTok}[1]{\textcolor[rgb]{0.56,0.13,0.00}{{#1}}}
    \newcommand{\DecValTok}[1]{\textcolor[rgb]{0.25,0.63,0.44}{{#1}}}
    \newcommand{\BaseNTok}[1]{\textcolor[rgb]{0.25,0.63,0.44}{{#1}}}
    \newcommand{\FloatTok}[1]{\textcolor[rgb]{0.25,0.63,0.44}{{#1}}}
    \newcommand{\CharTok}[1]{\textcolor[rgb]{0.25,0.44,0.63}{{#1}}}
    \newcommand{\StringTok}[1]{\textcolor[rgb]{0.25,0.44,0.63}{{#1}}}
    \newcommand{\CommentTok}[1]{\textcolor[rgb]{0.38,0.63,0.69}{\textit{{#1}}}}
    \newcommand{\OtherTok}[1]{\textcolor[rgb]{0.00,0.44,0.13}{{#1}}}
    \newcommand{\AlertTok}[1]{\textcolor[rgb]{1.00,0.00,0.00}{\textbf{{#1}}}}
    \newcommand{\FunctionTok}[1]{\textcolor[rgb]{0.02,0.16,0.49}{{#1}}}
    \newcommand{\RegionMarkerTok}[1]{{#1}}
    \newcommand{\ErrorTok}[1]{\textcolor[rgb]{1.00,0.00,0.00}{\textbf{{#1}}}}
    \newcommand{\NormalTok}[1]{{#1}}

    % Additional commands for more recent versions of Pandoc
    \newcommand{\ConstantTok}[1]{\textcolor[rgb]{0.53,0.00,0.00}{{#1}}}
    \newcommand{\SpecialCharTok}[1]{\textcolor[rgb]{0.25,0.44,0.63}{{#1}}}
    \newcommand{\VerbatimStringTok}[1]{\textcolor[rgb]{0.25,0.44,0.63}{{#1}}}
    \newcommand{\SpecialStringTok}[1]{\textcolor[rgb]{0.73,0.40,0.53}{{#1}}}
    \newcommand{\ImportTok}[1]{{#1}}
    \newcommand{\DocumentationTok}[1]{\textcolor[rgb]{0.73,0.13,0.13}{\textit{{#1}}}}
    \newcommand{\AnnotationTok}[1]{\textcolor[rgb]{0.38,0.63,0.69}{\textbf{\textit{{#1}}}}}
    \newcommand{\CommentVarTok}[1]{\textcolor[rgb]{0.38,0.63,0.69}{\textbf{\textit{{#1}}}}}
    \newcommand{\VariableTok}[1]{\textcolor[rgb]{0.10,0.09,0.49}{{#1}}}
    \newcommand{\ControlFlowTok}[1]{\textcolor[rgb]{0.00,0.44,0.13}{\textbf{{#1}}}}
    \newcommand{\OperatorTok}[1]{\textcolor[rgb]{0.40,0.40,0.40}{{#1}}}
    \newcommand{\BuiltInTok}[1]{{#1}}
    \newcommand{\ExtensionTok}[1]{{#1}}
    \newcommand{\PreprocessorTok}[1]{\textcolor[rgb]{0.74,0.48,0.00}{{#1}}}
    \newcommand{\AttributeTok}[1]{\textcolor[rgb]{0.49,0.56,0.16}{{#1}}}
    \newcommand{\InformationTok}[1]{\textcolor[rgb]{0.38,0.63,0.69}{\textbf{\textit{{#1}}}}}
    \newcommand{\WarningTok}[1]{\textcolor[rgb]{0.38,0.63,0.69}{\textbf{\textit{{#1}}}}}
    \makeatletter
    \newsavebox\pandoc@box
    \newcommand*\pandocbounded[1]{%
      \sbox\pandoc@box{#1}%
      % scaling factors for width and height
      \Gscale@div\@tempa\textheight{\dimexpr\ht\pandoc@box+\dp\pandoc@box\relax}%
      \Gscale@div\@tempb\linewidth{\wd\pandoc@box}%
      % select the smaller of both
      \ifdim\@tempb\p@<\@tempa\p@
        \let\@tempa\@tempb
      \fi
      % scaling accordingly (\@tempa < 1)
      \ifdim\@tempa\p@<\p@
        \scalebox{\@tempa}{\usebox\pandoc@box}%
      % scaling not needed, use as it is
      \else
        \usebox{\pandoc@box}%
      \fi
    }
    \makeatother

    % Define a nice break command that doesn't care if a line doesn't already
    % exist.
    \def\br{\hspace*{\fill} \\* }
    % Math Jax compatibility definitions
    \def\gt{>}
    \def\lt{<}
    \let\Oldtex\TeX
    \let\Oldlatex\LaTeX
    \renewcommand{\TeX}{\textrm{\Oldtex}}
    \renewcommand{\LaTeX}{\textrm{\Oldlatex}}
    % Document parameters
    % Document title
    \title{Homework 1\\[0.5ex]\large Computer Simulation 1}
    \author{Andrea Pokraj\v{c}i\'{c}\\
    \normalsize\href{mailto:andrea.pokrajcic@studserv.uni-leipzig.de}{andrea.pokrajcic@studserv.uni-leipzig.de}\\
    \normalsize Matriculation number: 3793956}
    \date{}
    
    
    
    
    
    
    
% Pygments definitions
\makeatletter
\def\PY@reset{\let\PY@it=\relax \let\PY@bf=\relax%
    \let\PY@ul=\relax \let\PY@tc=\relax%
    \let\PY@bc=\relax \let\PY@ff=\relax}
\def\PY@tok#1{\csname PY@tok@#1\endcsname}
\def\PY@toks#1+{\ifx\relax#1\empty\else%
    \PY@tok{#1}\expandafter\PY@toks\fi}
\def\PY@do#1{\PY@bc{\PY@tc{\PY@ul{%
    \PY@it{\PY@bf{\PY@ff{#1}}}}}}}
\def\PY#1#2{\PY@reset\PY@toks#1+\relax+\PY@do{#2}}

\@namedef{PY@tok@w}{\def\PY@tc##1{\textcolor[rgb]{0.73,0.73,0.73}{##1}}}
\@namedef{PY@tok@c}{\let\PY@it=\textit\def\PY@tc##1{\textcolor[rgb]{0.24,0.48,0.48}{##1}}}
\@namedef{PY@tok@cp}{\def\PY@tc##1{\textcolor[rgb]{0.61,0.40,0.00}{##1}}}
\@namedef{PY@tok@k}{\let\PY@bf=\textbf\def\PY@tc##1{\textcolor[rgb]{0.00,0.50,0.00}{##1}}}
\@namedef{PY@tok@kp}{\def\PY@tc##1{\textcolor[rgb]{0.00,0.50,0.00}{##1}}}
\@namedef{PY@tok@kt}{\def\PY@tc##1{\textcolor[rgb]{0.69,0.00,0.25}{##1}}}
\@namedef{PY@tok@o}{\def\PY@tc##1{\textcolor[rgb]{0.40,0.40,0.40}{##1}}}
\@namedef{PY@tok@ow}{\let\PY@bf=\textbf\def\PY@tc##1{\textcolor[rgb]{0.67,0.13,1.00}{##1}}}
\@namedef{PY@tok@nb}{\def\PY@tc##1{\textcolor[rgb]{0.00,0.50,0.00}{##1}}}
\@namedef{PY@tok@nf}{\def\PY@tc##1{\textcolor[rgb]{0.00,0.00,1.00}{##1}}}
\@namedef{PY@tok@nc}{\let\PY@bf=\textbf\def\PY@tc##1{\textcolor[rgb]{0.00,0.00,1.00}{##1}}}
\@namedef{PY@tok@nn}{\let\PY@bf=\textbf\def\PY@tc##1{\textcolor[rgb]{0.00,0.00,1.00}{##1}}}
\@namedef{PY@tok@ne}{\let\PY@bf=\textbf\def\PY@tc##1{\textcolor[rgb]{0.80,0.25,0.22}{##1}}}
\@namedef{PY@tok@nv}{\def\PY@tc##1{\textcolor[rgb]{0.10,0.09,0.49}{##1}}}
\@namedef{PY@tok@no}{\def\PY@tc##1{\textcolor[rgb]{0.53,0.00,0.00}{##1}}}
\@namedef{PY@tok@nl}{\def\PY@tc##1{\textcolor[rgb]{0.46,0.46,0.00}{##1}}}
\@namedef{PY@tok@ni}{\let\PY@bf=\textbf\def\PY@tc##1{\textcolor[rgb]{0.44,0.44,0.44}{##1}}}
\@namedef{PY@tok@na}{\def\PY@tc##1{\textcolor[rgb]{0.41,0.47,0.13}{##1}}}
\@namedef{PY@tok@nt}{\let\PY@bf=\textbf\def\PY@tc##1{\textcolor[rgb]{0.00,0.50,0.00}{##1}}}
\@namedef{PY@tok@nd}{\def\PY@tc##1{\textcolor[rgb]{0.67,0.13,1.00}{##1}}}
\@namedef{PY@tok@s}{\def\PY@tc##1{\textcolor[rgb]{0.73,0.13,0.13}{##1}}}
\@namedef{PY@tok@sd}{\let\PY@it=\textit\def\PY@tc##1{\textcolor[rgb]{0.73,0.13,0.13}{##1}}}
\@namedef{PY@tok@si}{\let\PY@bf=\textbf\def\PY@tc##1{\textcolor[rgb]{0.64,0.35,0.47}{##1}}}
\@namedef{PY@tok@se}{\let\PY@bf=\textbf\def\PY@tc##1{\textcolor[rgb]{0.67,0.36,0.12}{##1}}}
\@namedef{PY@tok@sr}{\def\PY@tc##1{\textcolor[rgb]{0.64,0.35,0.47}{##1}}}
\@namedef{PY@tok@ss}{\def\PY@tc##1{\textcolor[rgb]{0.10,0.09,0.49}{##1}}}
\@namedef{PY@tok@sx}{\def\PY@tc##1{\textcolor[rgb]{0.00,0.50,0.00}{##1}}}
\@namedef{PY@tok@m}{\def\PY@tc##1{\textcolor[rgb]{0.40,0.40,0.40}{##1}}}
\@namedef{PY@tok@gh}{\let\PY@bf=\textbf\def\PY@tc##1{\textcolor[rgb]{0.00,0.00,0.50}{##1}}}
\@namedef{PY@tok@gu}{\let\PY@bf=\textbf\def\PY@tc##1{\textcolor[rgb]{0.50,0.00,0.50}{##1}}}
\@namedef{PY@tok@gd}{\def\PY@tc##1{\textcolor[rgb]{0.63,0.00,0.00}{##1}}}
\@namedef{PY@tok@gi}{\def\PY@tc##1{\textcolor[rgb]{0.00,0.52,0.00}{##1}}}
\@namedef{PY@tok@gr}{\def\PY@tc##1{\textcolor[rgb]{0.89,0.00,0.00}{##1}}}
\@namedef{PY@tok@ge}{\let\PY@it=\textit}
\@namedef{PY@tok@gs}{\let\PY@bf=\textbf}
\@namedef{PY@tok@ges}{\let\PY@bf=\textbf\let\PY@it=\textit}
\@namedef{PY@tok@gp}{\let\PY@bf=\textbf\def\PY@tc##1{\textcolor[rgb]{0.00,0.00,0.50}{##1}}}
\@namedef{PY@tok@go}{\def\PY@tc##1{\textcolor[rgb]{0.44,0.44,0.44}{##1}}}
\@namedef{PY@tok@gt}{\def\PY@tc##1{\textcolor[rgb]{0.00,0.27,0.87}{##1}}}
\@namedef{PY@tok@err}{\def\PY@bc##1{{\setlength{\fboxsep}{\string -\fboxrule}\fcolorbox[rgb]{1.00,0.00,0.00}{1,1,1}{\strut ##1}}}}
\@namedef{PY@tok@kc}{\let\PY@bf=\textbf\def\PY@tc##1{\textcolor[rgb]{0.00,0.50,0.00}{##1}}}
\@namedef{PY@tok@kd}{\let\PY@bf=\textbf\def\PY@tc##1{\textcolor[rgb]{0.00,0.50,0.00}{##1}}}
\@namedef{PY@tok@kn}{\let\PY@bf=\textbf\def\PY@tc##1{\textcolor[rgb]{0.00,0.50,0.00}{##1}}}
\@namedef{PY@tok@kr}{\let\PY@bf=\textbf\def\PY@tc##1{\textcolor[rgb]{0.00,0.50,0.00}{##1}}}
\@namedef{PY@tok@bp}{\def\PY@tc##1{\textcolor[rgb]{0.00,0.50,0.00}{##1}}}
\@namedef{PY@tok@fm}{\def\PY@tc##1{\textcolor[rgb]{0.00,0.00,1.00}{##1}}}
\@namedef{PY@tok@vc}{\def\PY@tc##1{\textcolor[rgb]{0.10,0.09,0.49}{##1}}}
\@namedef{PY@tok@vg}{\def\PY@tc##1{\textcolor[rgb]{0.10,0.09,0.49}{##1}}}
\@namedef{PY@tok@vi}{\def\PY@tc##1{\textcolor[rgb]{0.10,0.09,0.49}{##1}}}
\@namedef{PY@tok@vm}{\def\PY@tc##1{\textcolor[rgb]{0.10,0.09,0.49}{##1}}}
\@namedef{PY@tok@sa}{\def\PY@tc##1{\textcolor[rgb]{0.73,0.13,0.13}{##1}}}
\@namedef{PY@tok@sb}{\def\PY@tc##1{\textcolor[rgb]{0.73,0.13,0.13}{##1}}}
\@namedef{PY@tok@sc}{\def\PY@tc##1{\textcolor[rgb]{0.73,0.13,0.13}{##1}}}
\@namedef{PY@tok@dl}{\def\PY@tc##1{\textcolor[rgb]{0.73,0.13,0.13}{##1}}}
\@namedef{PY@tok@s2}{\def\PY@tc##1{\textcolor[rgb]{0.73,0.13,0.13}{##1}}}
\@namedef{PY@tok@sh}{\def\PY@tc##1{\textcolor[rgb]{0.73,0.13,0.13}{##1}}}
\@namedef{PY@tok@s1}{\def\PY@tc##1{\textcolor[rgb]{0.73,0.13,0.13}{##1}}}
\@namedef{PY@tok@mb}{\def\PY@tc##1{\textcolor[rgb]{0.40,0.40,0.40}{##1}}}
\@namedef{PY@tok@mf}{\def\PY@tc##1{\textcolor[rgb]{0.40,0.40,0.40}{##1}}}
\@namedef{PY@tok@mh}{\def\PY@tc##1{\textcolor[rgb]{0.40,0.40,0.40}{##1}}}
\@namedef{PY@tok@mi}{\def\PY@tc##1{\textcolor[rgb]{0.40,0.40,0.40}{##1}}}
\@namedef{PY@tok@il}{\def\PY@tc##1{\textcolor[rgb]{0.40,0.40,0.40}{##1}}}
\@namedef{PY@tok@mo}{\def\PY@tc##1{\textcolor[rgb]{0.40,0.40,0.40}{##1}}}
\@namedef{PY@tok@ch}{\let\PY@it=\textit\def\PY@tc##1{\textcolor[rgb]{0.24,0.48,0.48}{##1}}}
\@namedef{PY@tok@cm}{\let\PY@it=\textit\def\PY@tc##1{\textcolor[rgb]{0.24,0.48,0.48}{##1}}}
\@namedef{PY@tok@cpf}{\let\PY@it=\textit\def\PY@tc##1{\textcolor[rgb]{0.24,0.48,0.48}{##1}}}
\@namedef{PY@tok@c1}{\let\PY@it=\textit\def\PY@tc##1{\textcolor[rgb]{0.24,0.48,0.48}{##1}}}
\@namedef{PY@tok@cs}{\let\PY@it=\textit\def\PY@tc##1{\textcolor[rgb]{0.24,0.48,0.48}{##1}}}

\def\PYZbs{\char`\\}
\def\PYZus{\char`\_}
\def\PYZob{\char`\{}
\def\PYZcb{\char`\}}
\def\PYZca{\char`\^}
\def\PYZam{\char`\&}
\def\PYZlt{\char`\<}
\def\PYZgt{\char`\>}
\def\PYZsh{\char`\#}
\def\PYZpc{\char`\%}
\def\PYZdl{\char`\$}
\def\PYZhy{\char`\-}
\def\PYZsq{\char`\'}
\def\PYZdq{\char`\"}
\def\PYZti{\char`\~}
% for compatibility with earlier versions
\def\PYZat{@}
\def\PYZlb{[}
\def\PYZrb{]}
\makeatother


    % For linebreaks inside Verbatim environment from package fancyvrb.
    \makeatletter
        \newbox\Wrappedcontinuationbox
        \newbox\Wrappedvisiblespacebox
        \newcommand*\Wrappedvisiblespace {\textcolor{red}{\textvisiblespace}}
        \newcommand*\Wrappedcontinuationsymbol {\textcolor{red}{\llap{\tiny$\m@th\hookrightarrow$}}}
        \newcommand*\Wrappedcontinuationindent {3ex }
        \newcommand*\Wrappedafterbreak {\kern\Wrappedcontinuationindent\copy\Wrappedcontinuationbox}
        % Take advantage of the already applied Pygments mark-up to insert
        % potential linebreaks for TeX processing.
        %        {, <, #, %, $, ' and ": go to next line.
        %        _, }, ^, &, >, - and ~: stay at end of broken line.
        % Use of \textquotesingle for straight quote.
        \newcommand*\Wrappedbreaksatspecials {%
            \def\PYGZus{\discretionary{\char`\_}{\Wrappedafterbreak}{\char`\_}}%
            \def\PYGZob{\discretionary{}{\Wrappedafterbreak\char`\{}{\char`\{}}%
            \def\PYGZcb{\discretionary{\char`\}}{\Wrappedafterbreak}{\char`\}}}%
            \def\PYGZca{\discretionary{\char`\^}{\Wrappedafterbreak}{\char`\^}}%
            \def\PYGZam{\discretionary{\char`\&}{\Wrappedafterbreak}{\char`\&}}%
            \def\PYGZlt{\discretionary{}{\Wrappedafterbreak\char`\<}{\char`\<}}%
            \def\PYGZgt{\discretionary{\char`\>}{\Wrappedafterbreak}{\char`\>}}%
            \def\PYGZsh{\discretionary{}{\Wrappedafterbreak\char`\#}{\char`\#}}%
            \def\PYGZpc{\discretionary{}{\Wrappedafterbreak\char`\%}{\char`\%}}%
            \def\PYGZdl{\discretionary{}{\Wrappedafterbreak\char`\$}{\char`\$}}%
            \def\PYGZhy{\discretionary{\char`\-}{\Wrappedafterbreak}{\char`\-}}%
            \def\PYGZsq{\discretionary{}{\Wrappedafterbreak\textquotesingle}{\textquotesingle}}%
            \def\PYGZdq{\discretionary{}{\Wrappedafterbreak\char`\"}{\char`\"}}%
            \def\PYGZti{\discretionary{\char`\~}{\Wrappedafterbreak}{\char`\~}}%
        }
        % Some characters . , ; ? ! / are not pygmentized.
        % This macro makes them "active" and they will insert potential linebreaks
        \newcommand*\Wrappedbreaksatpunct {%
            \lccode`\~`\.\lowercase{\def~}{\discretionary{\hbox{\char`\.}}{\Wrappedafterbreak}{\hbox{\char`\.}}}%
            \lccode`\~`\,\lowercase{\def~}{\discretionary{\hbox{\char`\,}}{\Wrappedafterbreak}{\hbox{\char`\,}}}%
            \lccode`\~`\;\lowercase{\def~}{\discretionary{\hbox{\char`\;}}{\Wrappedafterbreak}{\hbox{\char`\;}}}%
            \lccode`\~`\:\lowercase{\def~}{\discretionary{\hbox{\char`\:}}{\Wrappedafterbreak}{\hbox{\char`\:}}}%
            \lccode`\~`\?\lowercase{\def~}{\discretionary{\hbox{\char`\?}}{\Wrappedafterbreak}{\hbox{\char`\?}}}%
            \lccode`\~`\!\lowercase{\def~}{\discretionary{\hbox{\char`\!}}{\Wrappedafterbreak}{\hbox{\char`\!}}}%
            \lccode`\~`\/\lowercase{\def~}{\discretionary{\hbox{\char`\/}}{\Wrappedafterbreak}{\hbox{\char`\/}}}%
            \catcode`\.\active
            \catcode`\,\active
            \catcode`\;\active
            \catcode`\:\active
            \catcode`\?\active
            \catcode`\!\active
            \catcode`\/\active
            \lccode`\~`\~
        }
    \makeatother

    \let\OriginalVerbatim=\Verbatim
    \makeatletter
    \renewcommand{\Verbatim}[1][1]{%
        %\parskip\z@skip
        \sbox\Wrappedcontinuationbox {\Wrappedcontinuationsymbol}%
        \sbox\Wrappedvisiblespacebox {\FV@SetupFont\Wrappedvisiblespace}%
        \def\FancyVerbFormatLine ##1{\hsize\linewidth
            \vtop{\raggedright\hyphenpenalty\z@\exhyphenpenalty\z@
                \doublehyphendemerits\z@\finalhyphendemerits\z@
                \strut ##1\strut}%
        }%
        % If the linebreak is at a space, the latter will be displayed as visible
        % space at end of first line, and a continuation symbol starts next line.
        % Stretch/shrink are however usually zero for typewriter font.
        \def\FV@Space {%
            \nobreak\hskip\z@ plus\fontdimen3\font minus\fontdimen4\font
            \discretionary{\copy\Wrappedvisiblespacebox}{\Wrappedafterbreak}
            {\kern\fontdimen2\font}%
        }%

        % Allow breaks at special characters using \PYG... macros.
        \Wrappedbreaksatspecials
        % Breaks at punctuation characters . , ; ? ! and / need catcode=\active
        \OriginalVerbatim[#1,codes*=\Wrappedbreaksatpunct]%
    }
    \makeatother

    % Exact colors from NB
    \definecolor{incolor}{HTML}{303F9F}
    \definecolor{outcolor}{HTML}{D84315}
    \definecolor{cellborder}{HTML}{CFCFCF}
    \definecolor{cellbackground}{HTML}{F7F7F7}

    % prompt
    \makeatletter
    \newcommand{\boxspacing}{\kern\kvtcb@left@rule\kern\kvtcb@boxsep}
    \makeatother
    \newcommand{\prompt}[4]{
        {\ttfamily\llap{{\color{#2}[#3]:\hspace{3pt}#4}}\vspace{-\baselineskip}}
    }
    

    
    % Prevent overflowing lines due to hard-to-break entities
    \sloppy
    % Setup hyperref package
    \hypersetup{
      breaklinks=true,  % so long urls are correctly broken across lines
      colorlinks=true,
      urlcolor=urlcolor,
      linkcolor=linkcolor,
      citecolor=citecolor,
      }
    % Slightly bigger margins than the latex defaults
    
    \geometry{verbose,tmargin=1in,bmargin=1in,lmargin=1in,rmargin=1in}
    
    

\begin{document}
    
    \maketitle
    
    

    
    \section*{Problem 1: Statement}\label{problem-1-statement}

Calculate the exact solution of the one-dimensional Ising model with
\(N\) spins for free boundary conditions:

\begin{enumerate}
\def\labelenumi{\alph{enumi})}
\item
  Sketch the calculation of the partition function \(Z\) by successive
  summation of spins.
\item
  Calculate the internal energy
  \(U = \langle E\rangle = -\dfrac{\partial \ln Z}{\partial \beta}\) and
  the specific heat \(C = \dfrac{dU}{dT}\).
\item
  How do \(u = U/N\) and \(c = C/N\) behave in the thermodynamic limit
  \(N \to \infty\)?
\end{enumerate}

    \section*{Problem 1: Solution}\label{problem-1-solution}

    \subsection*{Part (a)}

Let us consider a one dimensional Ising model with \(N\) spins for free
boundary conditions. We can write the energy of such system as follows

\[
E=-J \sum_{i=1}^{N-1}s_i s_{i+1} \tag{1}
\]

where the sum is over all pairs of neighbouring spins \(i\) and \(i+1\)
and \(J\) is an exchange constant assumed to be positive. For algebraic
simplicity's sake, we can expand eq. (1) as follows

\[
E=-J(s_1s_2+s_2s_3+\cdots+s_{N-2}s_{N-1}+s_{N-1}s_N) \tag{2}
\]

Furthermore, by definition, we know that the partition function is given
by the following equation:

\[
Z=\sum_{\{s_i=\pm1\}} e^{\frac{-E}{k_B T}} \tag{3}
\]

Plugging in eq. (2) into eq. (3) we arrive at eq. (4)

\[
Z=\sum_{\{s_i=\pm1\}} e^{\frac{J}{k_B T}(s_1s_2+s_2s_3+\cdots+s_{N-2}s_{N-1}+s_{N-1}s_N)} \tag{4}
\]

This can be expanded over all states as follows:

\[
Z=\sum_{s_1=\pm1}\cdots\sum_{s_N=\pm1}e^{\frac{J}{k_B T}(s_1s_2+s_2s_3+\cdots+s_{N-2}s_{N-1}+s_{N-1}s_N)} \tag{5}
\]

Notice how the term \(s_N\) appears only in the last term. Therefore, we
can extract the sum accounting for the last term as shown in eq. (6)

\[
Z=\sum_{s_1=\pm1}\cdots\sum_{s_{N-1}=\pm1}e^{\frac{J}{k_B T}(s_1s_2+s_2s_3+\cdots+s_{N-2}s_{N-1})} \sum_{s_N=\pm1} e^{\frac{J}{k_B T}s_{N-1}s_N} \tag{6}
\]

Since \(s_N\) can only be \(+1\) or \(-1\), we calculate the last term
as:

\[
\sum_{s_N=\pm1} e^{\frac{J}{k_B T}s_{N-1}s_N} = e^{\frac{J}{k_B T}s_{N-1}} + e^{-\frac{J}{k_B T}s_{N-1}} \tag{7}
\]

Recalling the hyperbolic cosine definition given by eq. (8)

\[
\cosh(\theta)=\frac{e^{\theta}+e^{-\theta}}{2} \tag{8}
\]

One can rewrite eq. (7) in terms of the hyperbolic cosine. Since we know
that \(s_{N-1}=\pm1\) and since the hyperbolic cosine is an even
function we can further simplify the expression as follows

\[
\sum_{s_N=\pm1} e^{\frac{J}{k_B T}s_{N-1}s_N} = 2\cosh\left(\frac{Js_{N-1}}{k_B T}\right) = 2\cosh\left(\frac{J}{k_B T}\right) \tag{9}
\]

Plugging this back into eq. (6)

\[
Z=\sum_{s_1=\pm1}\cdots\sum_{s_{N-1}=\pm1}e^{\frac{J}{k_B T}(s_1s_2+s_2s_3+\cdots+s_{N-2}s_{N-1})}\, 2\cosh\left(\frac{J}{k_B T}\right) \tag{10}
\]

we realize this can be done recursively from the ``back'' of the
expression towards the front, i.e.:

\[
Z=\sum_{s_1=\pm1}\cdots\sum_{s_{N-2}=\pm1}e^{\frac{J}{k_B T}(s_1s_2+s_2s_3+\cdots+s_{N-3}s_{N-2})}\, 2\cosh\left(\frac{J}{k_B T}\right) 2\cosh\left(\frac{J}{k_B T}\right) \tag{11}
\]

\[
Z=\sum_{s_1=\pm1}\cdots\sum_{s_{N-3}=\pm1}e^{\frac{J}{k_B T}(s_1s_2+s_2s_3+\cdots+s_{N-4}s_{N-3})}\, 2\cosh\left(\frac{J}{k_B T}\right) 2\cosh\left(\frac{J}{k_B T}\right) 2\cosh\left(\frac{J}{k_B T}\right) \tag{12}
\]

and so on, \(N-1\) times, until you arrive at the following

\[
Z=\left[2\cosh\left(\frac{J}{k_B T}\right)\right]^{N-1}\sum_{s_1=\pm1}1 \tag{13}
\]

and since \(\sum_{s_1=\pm1}1=1+1=2\), finally we get:

\[
\boxed{Z=2^N\left[\cosh\left(\frac{J}{k_B T}\right)\right]^{N-1}} \tag{14}
\]

    \subsection*{Part (b)}

Now we are tasked with calculating the internal energy
\(U=\langle E\rangle=-\frac{\partial \ln Z}{\partial \beta}\) and the
specific heat \(C=\frac{dU}{dT}\).

First we calculate the internal energy by taking a partial derivative
with respect to \(\beta\), which is assumed to be equal to
\(\beta=\frac{1}{k_B T}\).

\[
\begin{aligned}
U &= -\frac{\partial \ln Z}{\partial \beta} = -\frac{\partial}{\partial \beta} \ln\left(2^N\left[\cosh(J\beta)\right]^{N-1}\right) \\
&= -\frac{\partial}{\partial \beta} \left(N\ln 2 + (N-1)\ln\left(\cosh(J\beta)\right)\right) \\
&= -\frac{\partial}{\partial \beta} (N-1) \ln\left(\cosh(J\beta)\right)
\end{aligned}
\]

by employing the chain rule we get

\[
= -(N-1)\frac{1}{\cosh(J\beta)}\sinh(J\beta)J
\]

furthermore using the definition of the hyperbolic tangent we conclude

\[
\boxed{U = -J(N-1)\tanh(J\beta)} \tag{15}
\]

or equivalently

\[
\boxed{U = -J(N-1)\tanh\left(\frac{J}{k_B T}\right)} \tag{16}
\]

If we wish to calculate the specific heat, we take a derivative of the
internal energy function with respect to temperature.

\[
\begin{aligned}
C &= \frac{dU}{dT} = \frac{d}{dT}\left(-J(N-1)\tanh\left(\frac{J}{k_B T}\right)\right) \\
&= -J(N-1)\frac{d}{dT}\tanh\left(\frac{J}{k_B T}\right)
\end{aligned}
\]

by knowing that
\(\frac{d \tanh(\theta)}{d\theta}=\operatorname{sech}^2(\theta)\) we
conclude

\[
\boxed{C = \frac{J^2}{k_B T^2}(N-1)\operatorname{sech}^2\left(\frac{J}{k_B T}\right)} \tag{17}
\]

    \subsection*{Part (c)}

Finally, we observe how \(u=U/N\) and \(c=C/N\) behave in the
thermodynamic limit \(N \to \infty\)

\[
u=\frac{-J(N-1)\tanh\left(\frac{J}{k_B T}\right)}{N}
\]

\[
c=\frac{\frac{J^2}{k_B T^2}(N-1)\operatorname{sech}^2\left(\frac{J}{k_B T}\right)}{N}
\]

taking the limit as \(N \to \infty\) for both expressions

\[
\begin{aligned}
&\lim_{N \to \infty} \frac{-J(N-1)\tanh\left(\frac{J}{k_B T}\right)}{N} \\
&= -J\tanh\left(\frac{J}{k_B T}\right) \lim_{N \to \infty} \frac{N-1}{N} \\
&= -J\tanh\left(\frac{J}{k_B T}\right) \lim_{N \to \infty} \left(1-\frac{1}{N}\right) \\
&\lim_{N \to \infty} u= \boxed{-J\tanh\left(\frac{J}{k_B T}\right)}
\end{aligned} \tag{18}
\]

\[
\begin{aligned}
&\lim_{N \to \infty} \frac{\frac{J^2}{k_B T^2}(N-1)\operatorname{sech}^2\left(\frac{J}{k_B T}\right)}{N} \\
&= \lim_{N \to \infty} \frac{J^2}{k_B T^2}\operatorname{sech}^2\left(\frac{J}{k_B T}\right)\frac{N-1}{N} \\
&= \frac{J^2}{k_B T^2}\operatorname{sech}^2\left(\frac{J}{k_B T}\right) \lim_{N \to \infty} \frac{N-1}{N} \\
&= \frac{J^2}{k_B T^2}\operatorname{sech}^2\left(\frac{J}{k_B T}\right) \lim_{N \to \infty} \left(1-\frac{1}{N}\right) \\
& \boxed{\lim_{N \to \infty} c=\frac{J^2}{k_B T^2}\operatorname{sech}^2\left(\frac{J}{k_B T}\right)}
\end{aligned} \tag{19}
\]

    \section*{Problem 2: Statement}\label{problem-2-statement}

\begin{enumerate}
\def\labelenumi{\alph{enumi})}
\item
  Plot \(u\) and \(c\) of exercise 1 as function of temperature \(T\)
  (e.g., with ``gnuplot'') for various finite \(N\)
  (\(= 20, 40, \dots, 100\)) and discuss the curves (for example: Does a
  maximum exist for \(c\) and, if yes, at which temperature?).
\item
  Study how the thermodynamic limit \(N \to \infty\) is approached?
\end{enumerate}

    \section*{Problem 2: Solution}\label{problem-2-solution}

    \subsection*{Part (a)}

We simplify expressions given by eq. (16) and (17)

\[
u(N,T)=-J\frac{(N-1)}{N}\tanh\left(\frac{J}{k_B T}\right) \tag{20}
\]

\[
c(N,T)=\frac{J^2}{k_B T^2}\frac{(N-1)}{N}\operatorname{sech}^2\left(\frac{J}{k_B T}\right) \tag{21}
\]

    Then we plot \(u\) as a function of \(T\) for several values of \(N\),
as demanded in the problem statement.

    \begin{tcolorbox}[breakable, size=fbox, boxrule=1pt, pad at break*=1mm,colback=cellbackground, colframe=cellborder]
\prompt{In}{incolor}{16}{\boxspacing}
\begin{Verbatim}[commandchars=\\\{\}]
\PY{k+kn}{import}\PY{+w}{ }\PY{n+nn}{numpy}\PY{+w}{ }\PY{k}{as}\PY{+w}{ }\PY{n+nn}{np}
\PY{k+kn}{from}\PY{+w}{ }\PY{n+nn}{matplotlib}\PY{+w}{ }\PY{k+kn}{import} \PY{n}{pyplot} \PY{k}{as} \PY{n}{plt}

\PY{n}{J} \PY{o}{=} \PY{l+m+mf}{1.0}
\PY{n}{k\PYZus{}b} \PY{o}{=} \PY{l+m+mf}{1.0}

\PY{n}{N\PYZus{}values} \PY{o}{=} \PY{p}{[}\PY{l+m+mi}{20}\PY{p}{,} \PY{l+m+mi}{40}\PY{p}{,} \PY{l+m+mi}{60}\PY{p}{,} \PY{l+m+mi}{80}\PY{p}{,} \PY{l+m+mi}{100}\PY{p}{]}

\PY{n}{T} \PY{o}{=} \PY{n}{np}\PY{o}{.}\PY{n}{linspace}\PY{p}{(}\PY{l+m+mf}{0.01}\PY{p}{,} \PY{l+m+mi}{50}\PY{p}{,} \PY{l+m+mi}{500}\PY{p}{)}

\PY{k}{def}\PY{+w}{ }\PY{n+nf}{u}\PY{p}{(}\PY{n}{N}\PY{p}{,}\PY{n}{T}\PY{p}{)}\PY{p}{:}
  \PY{k}{return} \PY{o}{\PYZhy{}}\PY{n}{J} \PY{o}{*} \PY{p}{(}\PY{n}{N}\PY{o}{\PYZhy{}}\PY{l+m+mi}{1}\PY{p}{)} \PY{o}{/} \PY{n}{N} \PY{o}{*} \PY{n}{np}\PY{o}{.}\PY{n}{tanh}\PY{p}{(}\PY{n}{J}\PY{o}{/}\PY{p}{(}\PY{n}{k\PYZus{}b}\PY{o}{*}\PY{n}{T}\PY{p}{)}\PY{p}{)}

\PY{k}{for} \PY{n}{N} \PY{o+ow}{in} \PY{n}{N\PYZus{}values}\PY{p}{:}
    \PY{n}{plt}\PY{o}{.}\PY{n}{plot}\PY{p}{(}\PY{n}{T}\PY{p}{,}\PY{n}{u}\PY{p}{(}\PY{n}{N}\PY{p}{,}\PY{n}{T}\PY{p}{)}\PY{p}{,} \PY{n}{label}\PY{o}{=}\PY{l+s+sa}{f}\PY{l+s+s2}{\PYZdq{}}\PY{l+s+s2}{N=}\PY{l+s+si}{\PYZob{}}\PY{n}{N}\PY{l+s+si}{\PYZcb{}}\PY{l+s+s2}{\PYZdq{}}\PY{p}{)}

\PY{n}{plt}\PY{o}{.}\PY{n}{xlabel}\PY{p}{(}\PY{l+s+s2}{\PYZdq{}}\PY{l+s+s2}{T}\PY{l+s+s2}{\PYZdq{}}\PY{p}{)}
\PY{n}{plt}\PY{o}{.}\PY{n}{ylabel}\PY{p}{(}\PY{l+s+s2}{\PYZdq{}}\PY{l+s+s2}{u(N,T)}\PY{l+s+s2}{\PYZdq{}}\PY{p}{)}
\PY{n}{plt}\PY{o}{.}\PY{n}{legend}\PY{p}{(}\PY{p}{)}
\PY{n}{plt}\PY{o}{.}\PY{n}{show}\PY{p}{(}\PY{p}{)}
\end{Verbatim}
\end{tcolorbox}

    \begin{center}
    \adjustimage{max size={0.9\linewidth}{0.9\paperheight}}{CompSim_HW1_files/CompSim_HW1_9_0.png}
    \end{center}
    { \hspace*{\fill} \\}
    
    Figure 1. Dependence of \(u(N,T)\) on \(T\), for \(N=20,40,60,80,100\)

    As we can see on Fig. 1, as \(T \to \infty\), \(u(N,T) \to 0\).
Additionally we see that as N is increased, the curves start off at
lower and lower \(u(N,T)\). More mathematically speaking as \(T \to 0\),
\(u(N,T) \to -J(N-1)/N\).

    Furthermore we plot \(c\) as a function of \(T\) for several values of
\(N\), as stated in the problem statement.

    \begin{tcolorbox}[breakable, size=fbox, boxrule=1pt, pad at break*=1mm,colback=cellbackground, colframe=cellborder]
\prompt{In}{incolor}{17}{\boxspacing}
\begin{Verbatim}[commandchars=\\\{\}]
\PY{k}{def}\PY{+w}{ }\PY{n+nf}{c}\PY{p}{(}\PY{n}{N}\PY{p}{,} \PY{n}{T}\PY{p}{)}\PY{p}{:}
    \PY{k}{return} \PY{n}{J}\PY{o}{*}\PY{o}{*}\PY{l+m+mi}{2}\PY{o}{/}\PY{p}{(}\PY{n}{k\PYZus{}b}\PY{o}{*}\PY{n}{T}\PY{o}{*}\PY{o}{*}\PY{l+m+mi}{2}\PY{p}{)} \PY{o}{*} \PY{p}{(}\PY{n}{N}\PY{o}{\PYZhy{}}\PY{l+m+mi}{1}\PY{p}{)}\PY{o}{/}\PY{n}{N} \PY{o}{*} \PY{l+m+mi}{1}\PY{o}{/}\PY{p}{(}\PY{n}{np}\PY{o}{.}\PY{n}{cosh}\PY{p}{(}\PY{n}{J}\PY{o}{/}\PY{p}{(}\PY{n}{k\PYZus{}b}\PY{o}{*}\PY{n}{T}\PY{p}{)}\PY{p}{)}\PY{o}{*}\PY{o}{*}\PY{l+m+mi}{2}\PY{p}{)}

\PY{k}{for} \PY{n}{N} \PY{o+ow}{in} \PY{n}{N\PYZus{}values}\PY{p}{:}
    \PY{n}{plt}\PY{o}{.}\PY{n}{plot}\PY{p}{(}\PY{n}{T}\PY{p}{,}\PY{n}{c}\PY{p}{(}\PY{n}{N}\PY{p}{,}\PY{n}{T}\PY{p}{)}\PY{p}{,} \PY{n}{label}\PY{o}{=}\PY{l+s+sa}{f}\PY{l+s+s2}{\PYZdq{}}\PY{l+s+s2}{N=}\PY{l+s+si}{\PYZob{}}\PY{n}{N}\PY{l+s+si}{\PYZcb{}}\PY{l+s+s2}{\PYZdq{}}\PY{p}{)}

\PY{n}{plt}\PY{o}{.}\PY{n}{xlabel}\PY{p}{(}\PY{l+s+s2}{\PYZdq{}}\PY{l+s+s2}{T}\PY{l+s+s2}{\PYZdq{}}\PY{p}{)}
\PY{n}{plt}\PY{o}{.}\PY{n}{ylabel}\PY{p}{(}\PY{l+s+s2}{\PYZdq{}}\PY{l+s+s2}{c(N,T)}\PY{l+s+s2}{\PYZdq{}}\PY{p}{)}
\PY{n}{plt}\PY{o}{.}\PY{n}{legend}\PY{p}{(}\PY{p}{)}
\PY{n}{plt}\PY{o}{.}\PY{n}{show}\PY{p}{(}\PY{p}{)}
\end{Verbatim}
\end{tcolorbox}

    \begin{center}
    \adjustimage{max size={0.9\linewidth}{0.9\paperheight}}{CompSim_HW1_files/CompSim_HW1_13_0.png}
    \end{center}
    { \hspace*{\fill} \\}
    
    Figure 2. Dependence of \(c(N,T)\) on \(T\), for \(N=20,40,60,80,100\)

    As we can see on Fig. 2, as \(T \to \infty\), \(c(N,T) \to 0\).
Additionally, there exists a clear maximum of \(c(N,T)\). \(c(N,T)\)
increases from zero, reaches a maximum near \(T=1\), and then decreases
towards zero.

Let us plot a plot with a smaller \(T\) bandwidth to identify at which
exact \(T\) value the maximum of \(c(N,T)\) appears . We use
\texttt{np.argmax} to find the largest values of \(c\) for each \(N\),
and then find the corresponding \(T\) at which it appears.

    \begin{tcolorbox}[breakable, size=fbox, boxrule=1pt, pad at break*=1mm,colback=cellbackground, colframe=cellborder]
\prompt{In}{incolor}{23}{\boxspacing}
\begin{Verbatim}[commandchars=\\\{\}]
\PY{n}{T\PYZus{}2} \PY{o}{=} \PY{n}{np}\PY{o}{.}\PY{n}{linspace}\PY{p}{(}\PY{l+m+mf}{0.01}\PY{p}{,} \PY{l+m+mi}{5}\PY{p}{,} \PY{l+m+mi}{500}\PY{p}{)}

\PY{k}{for} \PY{n}{N} \PY{o+ow}{in} \PY{n}{N\PYZus{}values}\PY{p}{:}
    \PY{n}{c\PYZus{}val} \PY{o}{=} \PY{n}{c}\PY{p}{(}\PY{n}{N}\PY{p}{,} \PY{n}{T\PYZus{}2}\PY{p}{)}
    \PY{n}{plt}\PY{o}{.}\PY{n}{plot}\PY{p}{(}\PY{n}{T\PYZus{}2}\PY{p}{,} \PY{n}{c\PYZus{}val}\PY{p}{,} \PY{n}{label}\PY{o}{=}\PY{l+s+sa}{f}\PY{l+s+s2}{\PYZdq{}}\PY{l+s+s2}{N=}\PY{l+s+si}{\PYZob{}}\PY{n}{N}\PY{l+s+si}{\PYZcb{}}\PY{l+s+s2}{\PYZdq{}}\PY{p}{)}

    \PY{n}{i\PYZus{}max} \PY{o}{=} \PY{n}{np}\PY{o}{.}\PY{n}{argmax}\PY{p}{(}\PY{n}{c\PYZus{}val}\PY{p}{)}
    \PY{n}{T\PYZus{}max} \PY{o}{=} \PY{n}{T\PYZus{}2}\PY{p}{[}\PY{n}{i\PYZus{}max}\PY{p}{]}

    \PY{n+nb}{print}\PY{p}{(}\PY{l+s+sa}{f}\PY{l+s+s2}{\PYZdq{}}\PY{l+s+s2}{N=}\PY{l+s+si}{\PYZob{}}\PY{n}{N}\PY{l+s+si}{\PYZcb{}}\PY{l+s+s2}{: T\PYZus{}max=}\PY{l+s+si}{\PYZob{}}\PY{n}{T\PYZus{}max}\PY{l+s+si}{:}\PY{l+s+s2}{.4f}\PY{l+s+si}{\PYZcb{}}\PY{l+s+s2}{\PYZdq{}}\PY{p}{)}

\PY{n}{plt}\PY{o}{.}\PY{n}{axvline}\PY{p}{(}\PY{l+m+mf}{0.83}\PY{p}{,} \PY{n}{linestyle}\PY{o}{=}\PY{l+s+s2}{\PYZdq{}}\PY{l+s+s2}{\PYZhy{}\PYZhy{}}\PY{l+s+s2}{\PYZdq{}}\PY{p}{,} \PY{n}{color}\PY{o}{=}\PY{l+s+s2}{\PYZdq{}}\PY{l+s+s2}{red}\PY{l+s+s2}{\PYZdq{}}\PY{p}{,} \PY{n}{label}\PY{o}{=}\PY{l+s+s2}{\PYZdq{}}\PY{l+s+s2}{T=0.83}\PY{l+s+s2}{\PYZdq{}}\PY{p}{)}
\PY{n}{plt}\PY{o}{.}\PY{n}{xlabel}\PY{p}{(}\PY{l+s+s2}{\PYZdq{}}\PY{l+s+s2}{T}\PY{l+s+s2}{\PYZdq{}}\PY{p}{)}
\PY{n}{plt}\PY{o}{.}\PY{n}{ylabel}\PY{p}{(}\PY{l+s+s2}{\PYZdq{}}\PY{l+s+s2}{c(N,T)}\PY{l+s+s2}{\PYZdq{}}\PY{p}{)}
\PY{n}{plt}\PY{o}{.}\PY{n}{legend}\PY{p}{(}\PY{p}{)}
\PY{n}{plt}\PY{o}{.}\PY{n}{show}\PY{p}{(}\PY{p}{)}
\end{Verbatim}
\end{tcolorbox}

    \begin{Verbatim}[commandchars=\\\{\}]
N=20: T\_max=0.8300
N=40: T\_max=0.8300
N=60: T\_max=0.8300
N=80: T\_max=0.8300
N=100: T\_max=0.8300
    \end{Verbatim}

    \begin{center}
    \adjustimage{max size={0.9\linewidth}{0.9\paperheight}}{CompSim_HW1_files/CompSim_HW1_16_1.png}
    \end{center}
    { \hspace*{\fill} \\}
    
    Figure 3. Dependence of \(c(N,T)\) on \(T\), for \(N=20,40,60,80,100\)
and \(T\le 5\)

    On Fig. 3 we can much more clearly see that the maximum of \(c(N,T)\)
exists at \(T\) around \(1\). Numerically we determine this number to be
\(T=0.83\). We also notice that increasing \(N\) changes the height of
the curves but not the temperature at which the heat capacity reaches
its maximum.

    \subsection*{Part (b)}

Now we wish to study how the thermodynamic limit of \(N \to \infty\)
behaves for both \(u(N,T)\) and \(c(N,T)\).

Let us remember eq. (18) and (19) derived in the previous task.

\[\lim_{N \to \infty} u= -J\tanh\left(\frac{J}{k_B T}\right)\]

\[\lim_{N \to \infty} c=\frac{J^2}{k_B T^2}\operatorname{sech}^2\left(\frac{J}{k_B T}\right)\]

It is useful to notice that \(u\) and \(c\) themselves given by eq. (20)
and (21) are nothing more than

\[u(N,T)=\frac{N-1}{N} \lim_{N \to \infty} u\]

\[c(N,T)=\frac{N-1}{N} \lim_{N \to \infty} c\]

Since the limits in eq. (18) and (19) are independent of \(N\), now we
see that the thermodynamic limit is approached through the factor
\(\frac{N-1}{N}=1-\frac{1}{N}\). This means that the error from the
thermodynamic limit decreases proportionally with \(\frac{1}{N}\).


    % Add a bibliography block to the postdoc
    
    
    
\end{document}
